\section{Proof of Lemma \ref{lem rank one}}\label{SM lem proof}
 
\begin{lemma} [Lemma \ref{lem rank one} restated]Suppose $\Sc$ is drawn from an LGP$(\sigma,\bSi,\bt_\st)$ as in Def.~\ref{def LGP} where $\text{rank}(\bSi)=1$ with $\bSi=\bla\bla^T$ for $\bla\in\R^p$. Define $\zeta=\Tb_s(\bla)^2/\tn{\bla}^2$. For magnitude and Hessian pruning ($P\in\{M,H\}$) and the associated retraining, we have the following excess risks with respect to $\bt^\st$%population minima
\begin{align}%\underbrace{\frac{\zeta^2\sigma^2}{p-2}}_{\text{Variance Term}}
&\E_{\Sc}[\Lc(\bth_s^P)]-\Lc(\bt^\st)={\frac{\zeta^2\sigma^2}{n-2}}+\underbrace{(1-\zeta)^2(\bla^T\bt^\st)^2}_{\text{Error due to bias}}\\
&\E_{\Sc}[\Lc(\bth_s^{RT})]-\Lc(\bt^\st)={\sigma^2}/({n-2}).
\end{align}
\end{lemma}
\begin{proof}
\noindent \textbf{Retraining analysis:} We claim that for any feature set $\Delta$ with $\bla_{\Delta}\neq 0$, the test risk of $\bth(\Delta)$ is exactly identical. Secondly, pruning is guaranteed to pick a  nonzero support satisfying $\bla_{\Delta}\neq 0$ \footnote{This is because $\bth=c\bla$ for some scalar $c\neq 0$ as $\bth$ lies in the row space of $\X$. Then, Hessian/Magnitude-pruning would pick a nonzero support of $\bth$ which corresponds to the nonzero support of $\bla$.}. Thus, as described next, retraining always achieves a fixed risk. Set $c^\st=\bla^T\bt^\st$. By definition, each input example $\x_i$ has the form $\x_i=g_i\bla$ and $y_i=g_ic^\st+\sigma z_i$. Set $\g=[g_1~\dots~g_n]^T$ and $\bar{\g}=\g/\tn{\g}$. Thus, we have $\X=\g\bla^T$ and $\y=\g\bla^T\bt^\st+\sigma\z$. Decompose $\z=\bar{\z}+\bar{\g}^T\z\bar{\g}$ where $\bar{\z}$ is orthogonal to $\g$. When solving the regression of $\Delta$, we have that 
\[
\X_\Delta=\g\blad^T,~\y=c^\st\g+\sigma(\bar{\z}+\bar{\g}^T\z\bar{\g})
\]
The least-squares solution has the form $\bth=\bth(\Delta)=\hat{c}\blad/\tn{\blad}^2$ where
\begin{align}
&\hat{c}=\arg\min_{c}\tn{(c^\st-c)\g+\sigma\z}\implies\hat{c}=c^\st+\sigma\gamma\label{beta}.
\end{align}
where $\gamma=\frac{\bar{\g}^T\z}{\tn{\g}}$. Observe that $\sqrt{p}\gamma$ has Student's t-distribution with $p$ degrees of freedom. Set $h$, $\epsilon\distas \Nn(0,1)$. Now, observe that a fresh test sample with $y=\x^T\bt^\st+\sigma \epsilon$ with $\x=g\bla$, the test error obeys
\begin{align}
\Lc(\bth(\Delta))&=\E(y-\x_{\Delta}^T\bth(\Delta))^2\\
&=\mathbb{E}[((c^\st-\hat{c})g+\sigma\epsilon)^2]\\
&=(\gamma^2+1)\sigma^2
\end{align}
Now, observe that the minimum risk is obviously $\Lc(\bt^\st)=\sigma^2$. Thus, the excess retraining risk becomes
\[
\Lc(\bth(\Delta))-\Lc(\bt^\st)=\gamma^2\sigma^2.
\]
regardless of choice of $\Delta$. Finally, averaging this risk over $\Sc$ returns $\E[\sigma^2\gamma^2]=\sigma^2/(n-2)$. Thus retraining risk has the fixed excess risk same as the one advertised in Lemma \ref{lem rank one}.

\noindent \textbf{Pruning analysis:} For pruning setting $\Delta=[p]$ above, we have that $\bth=\hat{c}\bla/\tn{\bla}^2$. This means that, for both Magnitude and Hessian pruning\footnote{They yield the same result since diagonal covariance is proportional to $\bla$ in magnitude.}, pruned vector takes the form $\bth_s=\hat{c}\Tb_s(\bla)/\tn{\bla}^2$. Using the fact that $\bla^T\Tb_s(\bla)=\tn{\Tb_s(\bla)}^2=\zeta\tn{\bla}^2$, we find
\begin{align}
\Lc(\bth_s)-\Lc(\bt^\st)&=\mathbb{E}[(c^\st- \hat{c}\frac{\bla^T\Tb_s(\bla)}{\tn{\bla}^2})g+\sigma\epsilon)^2]-\sigma^2\\
&=\mathbb{E}[(c^\st- \hat{c}\zeta)g+\sigma\epsilon)^2]-\sigma^2\\
&=\E[(((1-\zeta)c^\st-\zeta\sigma\gamma)g+\sigma\epsilon)^2]-\sigma^2\\
&=((1-\zeta)c^\st-\zeta\sigma\gamma)^2.
\end{align}
Finally, using zero-mean $\gamma$, we find
\[
\E_{\Sc}[\Lc(\bth_s)]-\Lc(\bt^\st)=\E_{\Sc}[((1-\zeta)c^\st-\zeta\sigma\gamma)^2]=(1-\zeta)^2{c^\st}^2+\frac{\zeta^2\sigma^2}{n-2},
\]
which concludes the proof after observing ${c^\st}^2=(\bla^T\bt^\st)^2$. Here, we call $(1-\zeta)^2{c^\st}^2$ ``the error due to bias''. The reason is that the predictable signal in the data is the noiseless component $\x^T\bt^\st$. Pruning leads to an error in this predictable component by resulting in a biased estimate of the label (when conditioned on the random variable $g$ which controls the signal).
\cmt{
Now let $|\lambda_{k_1}|\geq|\lambda_{k_2}|\geq...\geq|\lambda_{k_p}|$. Assume we solved ERM to get  $\hat{c}$ and $\bth=\hat{c}\bla$. Prune the trained model $\bth$ to $s$-sparse by keeping the largest entries. Set $\ub=\Tb_s(\bla)$ and $\zeta=\tn{\ub}^2/\tn{\bla}^2$. We get $\bth_s=\hat{c}\ub$.
\begin{align}
    \Lc(\tbh)&=\mathbb{E}[((\bla^T\bt-\ub^T\tbh)h+\sigma\epsilon)^2]\\
    &=\mathbb{E}[((\bla^T\bt-\sum_{i=1}^s\lambda_{k_i}^2\hat{c})h+\sigma\epsilon)^2]\\
    &=[(1-\zeta)\bla^T\bt-\zeta\sigma\gamma]^2+\sigma^2
\end{align}
\textbf{Case 1: }
Now, assume that we have already known the optimal entries and do pruning before training. Then the pruned model can be seen as pruning the inputs and then putting it through a non-sparse $s$-feature model. Let $\lambda_{t_1}\bt_{t_1}\geq\lambda_{t_2}\bt_{t_2}\geq...\geq\lambda_{t_s}\bt_{t_s}$, $\w=[\lambda_{t_1}~\lambda_{t_2}~...~\lambda_{t_s}]$ and $\tb=[\bt_{t_1}~...~\bt_{t_s}]$. Set data sample $\x_i'\distas \Nn(0,\bSi')\in\R^s$ where are $\bSi'=\w\w^T$. Following what done in \ref{beta}, similarly
\[
\tbh'=\theta\w\quad\text{where}\quad\theta=\frac{1}{\sum_{i=1}^s\lambda_{t_i}^2}(\w^T\tb+\frac{\sigma}{\tn{\g}^2}\g^T\z)
\]
\begin{align}
    \Lc(\tbh')&=\mathbb{E}[((\bla^T\bt-\sum_{i=1}^s\lambda_{t_i}^2\theta)h+\sigma\epsilon)^2]\\
    &=[(\bla^T\bt-\w^T\tb)-\sigma\gamma]^2+\sigma^2
\end{align}
\textbf{Case 2: }
Now retrain with the pruned entries $\lambda_{k_!}$, ..., $\lambda_{k_s}$. Then the pruned model can be seen as pruning the inputs and then putting it through a non-sparse $s$-feature model. Let $\tb=[\bt_{k_1}~...~\bt_{k_s}]$. Set data sample $\x_i'\distas \Nn(0,\bSi')\in\R^s$ where are $\bSi'=\ub\ub^T$. Following what done in \ref{beta}, similarly
\[
\tbh'=\theta\vct u\quad\text{where}\quad\theta=\frac{1}{\sum_{i=1}^s\lambda_{k_i}^2}(\vct u^T\tb+\frac{\sigma}{\tn{\g}^2}\g^T\z)
\]
\begin{align}
    \Lc(\tbh')&=\mathbb{E}[((\bla^T\bt-\sum_{i=1}^s\lambda_{k_i}^2\theta)h+\sigma\epsilon)^2]\\
    &=[(\bla^T\bt-\vct u^T\tb)-\sigma\gamma]^2+\sigma^2
\end{align}
\textbf{It is obvious to know that loss in case 1 is always smaller than it in case 2. Then for both case, }

Set $\bt^T\bSi\bt=\alpha^2$ and $\tb^T\bSi'\tb=\rho^2$. We can write $\Lc(\tbh)$ and $\Lc(\tbh')$ as
\[
\Lc(\tbh)=((1-\zeta)\alpha-\zeta\sigma\gamma)^2+\sigma^2\iff\Lc(\tbh')=(\alpha-\rho-\sigma\gamma)^2+\sigma^2
\]
Now consider a special case which satisfies $\Lc(\tbh)<\Lc(\tbh')$.
\[
0\leq(1-\zeta)\alpha-\zeta\sigma\gamma<\alpha-\rho-\sigma\gamma
\]
\[
\frac{\rho+\sigma\gamma}{\alpha+\sigma\gamma}<\zeta\leq\frac{\alpha}{\alpha+\sigma\gamma}
\]}
\end{proof}